\scanpage{96}

\uline{Prop.} Soit $K$ un corps de \unsure{nb} fini sur $Q$, $L$ une extension quadratique de $K$, $p_0$ un idéal premier de $K$, \deleted{$p,p'$} qui se décompose dans $L$ en deux idéaux premiers distincts $p,p'$. Alors il existe un $x\in L^*$ \deleted{dont le diviseur \unsure{doit être}} \deleted{tel que} dont le diviseur soit $h(p-p')$, \inserted{où $h\in N^+$,} et tel que $x\,{}^\sigma x=1$, où $\sigma$ est l'autom. non trivial de $L/K$.

\margin{Plus précisément à torsion près, les diviseurs des \deleted{\unsure{nb}} éléments de $L$ tels que $N_{L/K}x=1$, \deleted{\unsure{$N_{L/K}$}} sont les diviseurs de $L$ de norme $1$!}

En effet, \deleted{comme le groupe des} si $h$ est l'ordre du groupe de classes \inserted{de div.} de $L$, il existe un $y$ de diviseur $h(p-p')$. Alors $y\,{}^\sigma y$ est un élément de $K$ dont le diviseur est évidemment nul, donc c'est une unité $u$.
Donc si $z=y^2$, alors \deleted{$\unsure{x\bar x}\quad\unsure{x\,{}^\sigma x}$} $z\cdot{}^\sigma z=(y\cdot{}^\sigma y)^2=u^2$, donc si $x=\dfrac{z}{u}=\dfrac{y^2}{Ny}$, on aura
\[
\operatorname{div}x=\operatorname{div}z=2\operatorname{div}y=2h(p-p'),\qquad\text{et}\qquad x\,{}^\sigma x=(z\,{}^\sigma z)/u^2=1.
\]

\uline{N.B.} Supposons que $K$ est tot. réel \deleted{\unsure{galoisienne de $Q$}}, \deleted{et $L$ une \unsure{ext.}} \inserted{$L$ \deleted{\unsure{ext.}} qu. de $K$; supposée \uline{galoisienne} sur $Q$,} \deleted{\unsure{et $L$}} tot. imaginaire. Soit $p$ un nb premier, et $p_0$ idéal de $K$ au dessus de $p$, supposons qu'il n'y ait qu'un \uline{seul idéal} \inserted{premier} $p$ \deleted{\unsure{de}} de $L$ sur $K$. Soit $\lambda$ un \uline{entier} de $L$ tel que $\lambda\bar\lambda=p^i$. Alors $\lambda^2\bar\lambda^2=p^{2i}$ donc si $\mu=\lambda^2/p^i=\dfrac{\lambda^2}{N_{L/K}\lambda}$ on aura $\mu\bar\mu=1$, \deleted{et $\mu$ est une unité} qui implique que $\mu$ est une unité en $p$, donc ($K$ est galoisien) pour $\forall$ idéal premier au dessus de $p$ \unsure{par $Q$-conjugaison|puisque $\mu=\bar\mu^{-1}$} (l'hyp. sur $p_0$ \unsure{se transpose aux $p_0$ conj.}) donc $\mu$ est une unité. (NB. $\lambda$ \unsure{aux $p$|aussi} est une unité \unsure{exc. $p$}.) donc une racine de l'unité. Donc $\lambda=p^{i/2}u$, $u$ racine de l'unité!

\mspage{95}
